\chapter{The Exact Complementary-Cancellation Theorem}
\label{ch:cancellation}

\section{Scalar readout of a two-branch state}
Consider the coherent scalar readout
\begin{equation}
\mathcal A(\sigma,\phi)
=\sqrt{\sigma}+e^{i\phi}\sqrt{1-\sigma},
\qquad 0\leq\sigma\leq1.
\label{eq:amp}
\end{equation}
This is the overlap of the state \eqref{eq:qubit-state} with the unnormalized symmetric readout bra $\bra{0}+\bra{1}$. The associated intensity is
\begin{align}
\mathcal I(\sigma,\phi)
&=|\mathcal A(\sigma,\phi)|^2\\
&=\sigma+(1-\sigma)+2\sqrt{\sigma(1-\sigma)}\cos\phi\\
&=1+2\sqrt{\sigma(1-\sigma)}\cos\phi.
\label{eq:intensity}
\end{align}

\begin{theorem}[Unique exact cancellation]
\label{thm:cancellation}
For $0\leq\sigma\leq1$ and real $\phi$,
\begin{equation}
\mathcal A(\sigma,\phi)=0
\end{equation}
if and only if
\begin{equation}
\sigma=\frac12,
\qquad
\phi=(2k+1)\pi,
\quad k\in\mathbb Z.
\end{equation}
\end{theorem}

\begin{proof}
Suppose $\mathcal A=0$. Then
\begin{equation}
\sqrt{\sigma}=-e^{i\phi}\sqrt{1-\sigma}.
\end{equation}
Taking absolute values gives $\sqrt{\sigma}=\sqrt{1-\sigma}$, hence $\sigma=1/2$. Substituting this back yields $1+e^{i\phi}=0$, so $e^{i\phi}=-1$ and $\phi=(2k+1)\pi$. Conversely, those conditions clearly give $\mathcal A=0$.
\end{proof}

This theorem is the cleanest rigorous core of the information-spinor hypothesis. Exact zero formation in a normalized complementary pair requires both equal weights and phase opposition. Neither condition alone is sufficient.

\section{Alternative proof from the intensity bound}
Since
\begin{equation}
0\leq\sqrt{\sigma(1-\sigma)}\leq\frac12,
\end{equation}
with equality only at $\sigma=1/2$, equation \eqref{eq:intensity} implies
\begin{equation}
\mathcal I(\sigma,\phi)
\geq 1-2\sqrt{\sigma(1-\sigma)}\geq0.
\end{equation}
Equality in the first inequality requires $\cos\phi=-1$; equality in the second requires $\sigma=1/2$. This again proves Theorem \ref{thm:cancellation}.

\section{Factorization of the residual at phase pi}
At $\phi=\pi$,
\begin{equation}
\mathcal I(\sigma,\pi)
=\left(\sqrt{\sigma}-\sqrt{1-\sigma}\right)^2.
\end{equation}
Rationalizing gives
\begin{equation}
\mathcal I(\sigma,\pi)
=\frac{(2\sigma-1)^2}{\left(\sqrt{\sigma}+\sqrt{1-\sigma}\right)^2}.
\label{eq:intensity-delta}
\end{equation}
Hence the cancellation residual is an even, nonnegative function of the centered coordinate $\delta=\sigma-1/2$ and vanishes quadratically at the balance axis.

Near $\delta=0$,
\begin{equation}
\mathcal I\!\left(\frac12+\delta,\pi\right)
=2\delta^2+2\delta^4+O(\delta^6).
\end{equation}
The leading coefficient matches the leading Shannon entropy deficit derived in Chapter \ref{ch:shannon}:
\begin{equation}
\ln2-\Hbin\!\left(\frac12+\delta\right)
=2\delta^2+O(\delta^4).
\end{equation}
Thus, infinitesimally, informational imbalance and destructive-interference failure have the same quadratic cost.

\begin{proposition}[Local equivalence of balance costs]
As $\delta\to0$,
\begin{equation}
\frac{\mathcal I(1/2+\delta,\pi)}{\ln2-\Hbin(1/2+\delta)}\longrightarrow1.
\end{equation}
\end{proposition}

This local equality is stronger than the observation that both quantities have minima at $1/2$. It shows that, to second order, the Shannon information deficit is numerically identical to the loss of perfect phase cancellation in the normalized two-branch state.

\section{Geometric interpretation}
At $\sigma=1/2$, the state is
\begin{equation}
\ket{\psi(1/2,\phi)}
=\frac{1}{\sqrt2}\left(\ket0+e^{i\phi}\ket1\right).
\end{equation}
The points $\phi=0$ and $\phi=\pi$ are antipodal equatorial states. The cancellation readout at $\phi=\pi$ is orthogonal to the symmetric readout direction. Thus the zero is represented by a projective orthogonality event at the same equator that carries the $-1$ Berry holonomy.

\section{Generalized readout weights}
For completeness, consider
\begin{equation}
\mathcal A_{a,b}(\sigma,\phi)
=a\sqrt{\sigma}+b e^{i\phi}\sqrt{1-\sigma},
\end{equation}
with positive real $a,b$. Exact cancellation requires
\begin{equation}
\frac{\sigma}{1-\sigma}=\frac{b^2}{a^2},
\end{equation}
so
\begin{equation}
\sigma=\frac{b^2}{a^2+b^2}.
\end{equation}
Therefore $1/2$ is selected only by a symmetric readout $a=b$. The symmetry of the zeta functional equation is what motivates equal branch normalization. A canonical construction must derive, not merely choose, that symmetry.

\section{What this theorem does and does not establish}
The theorem proves the implication
\begin{equation}
\text{normalized complementary cancellation}
\Longrightarrow
\sigma=\frac12.
\end{equation}
It does not prove
\begin{equation}
\zeta(s)=0
\Longrightarrow
\text{normalized complementary cancellation}.
\end{equation}
The second implication is the open representation problem. Every later claim is organized around this logical boundary.

\begin{figure}[ht]
\centering
\includegraphics[width=0.76\textwidth]{destructive_interference.pdf}
\caption{At relative phase $\pi$, unequal complementary amplitudes leave a positive residual. Exact cancellation occurs only at $\sigma=1/2$.}
\label{fig:cancellation}
\end{figure}
